\chapter{Probabilités finies et dénombrement}
\section*{Exercice 21.1 -- Construire une loi}
Une variable prend les valeurs $1,2,3,4$ avec des probabilités proportionnelles à ces nombres. Déterminer la loi et la probabilité d'un résultat pair.

\section*{Exercice 21.2 -- Deux dés}
On lance deux dés équilibrés indépendants. Calculer la probabilité que la somme soit $7$, puis qu'au moins un dé donne $6$. Justifier l'univers utilisé.

\section*{Exercice 21.3 -- Réunion et intersection}
On sait $P(A)=0{,}4$, $P(B)=0{,}5$ et $P(A\cap B)=0{,}2$. Calculer $P(A\cup B)$ et la probabilité de n'avoir ni A ni B. Démontrer la formule employée en séparant trois régions disjointes.

\section*{Exercice 21.4 -- Remettre ou ne pas remettre}
Une urne contient trois boules rouges et deux bleues, indiscernables au toucher. On en tire deux. Calculer la probabilité de deux rouges avec remise, puis sans remise.

\section*{Exercice 21.5 -- Un modèle est une hypothèse}
Une pièce a donné $63$ faces en $100$ lancers. Peut-on conclure que sa probabilité de face vaut exactement $0{,}63$ ? Qu'apporterait une simulation de pièces équilibrées ?

\section*{Exercice 21.6 -- Compter un complément}
Un code comporte trois chiffres, chacun choisi indépendamment et uniformément parmi $0,\ldots,9$. Les répétitions et le zéro initial sont autorisés. Calculer la probabilité d'au moins une répétition.
